\hspace{3mm}

\centerline{\bf \huge 9 класс}
\centerline{\normalsize{\it Работа рассчитана на \bf{180} минут}}

\hspace{3mm}

\problem{9.1} Два друга записали на доске следующие суммы: Адьян $2023^{2023} + 2021^{2021}$, а Церен $2023^{2021} + 2021^{2023}.$ Определите, у кого из друзей получилась большая сумма?\

\hspace{3mm}

Ответ: У Адьяна.

Решение:\\
Вычтем из суммы Адьяна сумму Церена:\\
$2023^{2023} + 2021^{2021} - (2023^{2021} + 2021^{2023}) = 2023^{2023} - 2023^{2021} - (2021^{2023} - 2021^{2021}) = 2023^{2021}(2023^2 - 1) - 2021^{2021}(2021^2 - 1).$\\
Заметим, что $2023^{2021} > 2021^{2021} > 0 $ и  $2023^2 - 1 > 2021^2 - 1 > 0.$\\
Следовательно, уменьшаемое $(2023^{2021}(2023^2 - 1))$ больше вычитаемого\\
$(2021^{2021}(2021^2 - 1))$, то есть разность положительна. Значит, у Адьяна сумма больше.

\hspace{3mm}

Критерии:\\
Любое правильное решение с верным обоснованием: \textbf{7 баллов.}\\
Идеи, продвигающие ход решения, но недостаточные для получения верного ответа: \textbf{2 балла.}\\
Неверное решение или идеи, не продвигающие ход решения, такие как "видно, что одно из чисел больше другого" или "очевидно, что первое число больше...": \textbf{0 баллов.}\\
Верный или неверный ответ без обоснования: \textbf{0 баллов.}

\hspace{3mm}

\problem{9.2} На городской этап олимпиады по математике пришли 2023 ученика из разных школ Элисты, при этом у любых двух участников олимпиады разное количество знакомых среди других участников. Найдётся ли ученик, который будет знаком со всеми другими участниками городской олимпиады по математике? 

\hspace{3mm}

Ответ: Нет, не найдётся.

Решение:\\
Предположим, что нашёлся ученик, который знаком со всеми, и назовём его учеником X. Заметим, что если у любых двух разное количество знакомых, то у каждого разное количество знакомых среди других участников олимпиады. Всего 2023 всевозможных 

\newpage

\hspace{3mm}

\noindentвариантов количества знакомых у одного человека {0, 1, 2, …, 2022}. Учеников – 2023, при этом у каждого разное количество знакомых. Следовательно, будут использованы все 2023 варианта количества знакомых, в том числе и вариант с 0 знакомых. То есть найдётся ученик, который ни с кем не знаком, и в том числе он не будет знаком с учеником X, но по предположению ученик X знаком со всеми. Противоречие. Следовательно, наше предположение о том, что нашёлся такой ученик, который знаком со всеми, неверно. Значит, такой ученик не найдётся. 

\hspace{3mm}

Критерии:\\
Любое правильное решение с верным обоснованием: \textbf{7 баллов.}\\
Идеи, продвигающие ход решения, но недостаточные для получения верного ответа: \textbf{2 балла.}\\
Неверное решение или идеи, не продвигающие ход решения: \textbf{0 баллов.}\\
Верный или неверный ответ без обоснования: \textbf{0 баллов.}

\hspace{3mm}

\problem{9.3} Учитель выписал на доску все натуральные числа от 1 до 2023 и предложил Баиру каждую минуту выполнять следующую операцию до тех пор, пока у него не останется одно число на доске: стирать любые два числа $a$ и $b$ с доски и записывать вместо них одно число, равное $2022\cdot(ab)^{2023} + 1^{2023}$. Баир согласился и быстро справился с заданием. Какой остаток при делении на два даёт последнее оставшееся число на доске?

\hspace{3mm}

Ответ: 1.

Решение:\\
В любой момент времени на доске будут только натуральные числа, так как число $2022\cdot(ab)^{2023} + 1^{2023}$ является суммой двух натуральных чисел, а значит, оно натуральное.
Заметим, что после стирания любых двух чисел у нас всегда остается нечётное число, так как $2022\cdot(ab)^{2023}$ - число чётное при любых $a$ и $b$. Из-за того, что множитель $2022$ чётен, $1^{2023} = 1$ число нечётное, а сумма чётного и нечётного числа всегда является нечётной.
Значит, последнее число на доске - нечётное, и оно даёт остаток $1$ при делении на $2$.

\newpage

\hspace{3mm}

Критерии:\\
Приведено полное обоснованное решение: \textbf{7 баллов.}\\
Верное решение без обоснования, что в любой момент времени на доске будут только натуральные числа: \textbf{5 баллов.}\\
Идеи, продвигающие ход решения, но недостаточные для получения верного ответа: \textbf{2 балла.}\\
Неверное решение или идеи, не продвигающие ход решения: \textbf{0 баллов.}\\
Верный или неверный ответ без обоснования: \textbf{0 баллов.}


\problem{9.4} В параллелограмме $ABCD$ проведены диагонали $AC$ и $BD$, которые пересекаются в точке $O$. Точка $M$ - середина стороны $CD$. Точка $K$ делит диагональ $AC$ в отношении $AK$:$KC = 1:2$. Докажите, что прямые $MO$ и $DK$ пересекаются на середине стороны $AB$.

\hspace{3mm}

Решение:\\
Сперва докажем, что $MO$ попадает на середину $AB$. Действительно, $MO$ параллельна $AD$, и точка $O$ - это середина $BD$. Значит, продолжение $MO$ будет средней линией треугольника $ABD$ и пересечёт в середине $AB$. Теперь докажем, что $DK$ попадёт на середину $AB$. Точка K делит $AC$ в отношении $1:2$, но AC состоит из двух равных отрезков $AO$ и $OC$. Пусть $AK = x$. Тогда $AC = 3x$, $AO = 1,5x$. Стало быть, $AK/AO = \dfrac{x}{1,5x} = 2:3$. А значит $AK$ относится к $KO$ как $2:1$. Помимо этого мы знаем, что $AO$ - это медиана, а значит, что медиана делит отрезок $AO$ в отношении $2:1$, следовательно, $K$ - это точка пересечения медиан треугольника $ABD$. Значит, продолжение $DK$ - это тоже медиана, и она попадёт на середину $AB$.

\begin{figure}[h] 
    \includegraphics[scale=0.3]{Городская Олимпиада Элиста 2023/Решения/picture_9.4.jpg}
\end{figure}


\newpage

\hspace{3mm}

Критерии:\\
Приведено полное обоснованное решение: \textbf{7 баллов.}\\
Показано, что $DK$ попадает на середину $AB$: \textbf{5 баллов.}\\
Показано, что продолжение $MO$ попадает на середину $A$: \textbf{2 балла.}\\
Неверное решение или идеи, не продвигающие ход решения: \textbf{0 баллов.}


\hspace{3mm}

\problem{9.5 } Дан правильный 97-угольник. У Алтаны есть 10 банок с красками различных цветов. Она посчитала, сколькими способами может покрасить вершины этого 97-угольника своими красками, но те раскраски, которые переходят друг в друга поворотами, Алтана считает одинаковыми. Какое число у нее получилось? 
Указание: 97 - простое число

\hspace{3mm}

Ответ: $\dfrac{10^{97} - 10}{97} + 10.$

Решение:\\
Посчитаем количество раскрасок, которые состоят ровно из одного цвета. Их в точности 10, так как цветов у нас 10. Всего количество способов раскрасить $97$ вершин в $10$ цветов без учета поворотов - $10^{97}$. Вычтем из этого количества все одноцветные раскраски. Получим, что количество раскрасок, которые состоят минимум из двух цветов, равно $10^{97} - 10$. Рассмотрим любую такую раскраску. Минимальный угол, на который мы можем ее повернуть - это когда вершина переходит в соседнюю. Так как $97$ число простое, то, поворачивая такую раскраску на минимальный угол, мы каждый раз получаем "новую" \ раскраску, которой до этого не было. После $97$ таких поворотов мы придем в исходное положение. Получаем, что каждой не одноцветной раскраске соответствует $97$ ее поворотов. Но так как мы считаем одинаковыми те раскраски, которые переходят друг в друга поворотами, то нужно $10^{97} - 10$ поделить на $97$. Наконец, вернем все одноцветные раскраски, которые мы убрали, и получим ответ:\\

$\dfrac{10^{97} - 10}{97} + 10.$


\newpage

\hspace{3mm}

Критерии:\\
Полное верное решение: \textbf{7 баллов.}\\
Показано, что не одноцветные раскраски делятся на группы по 97, но в конце не добавлены одноцветные раскраски: \textbf{6 баллов.}\\
Посчитано, сколько раскрасок без учета поворотов: \textbf{2 балла.}\\
Посчитано количество одноцветных раскрасок: \textbf{1 балл.}\\
Неверное решение или идеи, не продвигающие ход решения: \textbf{0 баллов.}\\
Верный или неверный ответ без обоснования: \textbf{0 баллов
}